\documentclass[title=Monsieur, name=Baba, r=0.86, g=0.06, b=0.15]{../mrms}
\begin{document}

\coverpage{../img/m-baba.jpg}
\collectionpage
\titlepage

% On fait une série de tranches de vie
% Y'a une sortie au parc avec mme gobelin, m je suis pas fatigué et mme je veux pas dormir
% c'est sur les dérivées linguistiques des enfants
% Ça va prendre le trois de carreaux aussi. mMMMH
% Y'a l'halloween et le leadup à l'halloween

\txtpage{
    C'était un matin de fin d'été resplendissant. Monsieur Baba se réveilla d'un bond, et regarda
    son calendrier.

    C'était le samedi! Il s'empressa d'enfiler son linge, puis il appela ses amis.
}
\imgpage{./img/page-01.jpg}

\txtpage{
    - Vous avez vu? On est en congelé aujourd'hui! Mon travail est en congelé.
    On devrait aller au parc pour en profiter, annonça monsieur Baba au téléphone
    à monsieur Je-Suis-Pas-Fatigué, à madame Gobelin, et à madame Je-Veux-Pas-Dormir.
}
\imgpage{./img/page-02.jpg}

\txtpage{
    - S'il fait si froid, on veut vraiment aller dehors? Demanda madame Je-Veux-Pas-Dormir.

    Madame Gobelin, qui regardait son thermomètre accroché dehors, était perplexe. Il faisait doux.

    - Excellente idée, repris monsieur Je-Suis-Pas-Fatigué. Un beau matin de congelé comme ça au parc,
    c'est parfait.

    - Madame Gobelin gloussa un peu, et corrigea: Vous voulez dire qu'on est en congé?
}
\imgpage{./img/page-03.jpg}

\txtpage{
    Arrivés au parc, ils commencèrent à jouer au roi de la tour. Monsieur Je-Suis-Pas-Fatigué avait
    expliqué que c'est comme le roi de la montagne, sauf que le module est une tour, alors c'est
    différent.

    Monsieur Baba était bien d'accord, et madame Gobelin trouvait que c'était un peu long avant de
    commencer à grimper.
}
\imgpage{./img/page-04.jpg}

\txtpage{
    Monsieur Baba fut le premier en haut de la tour. Il était le roi maintenant.

    Madame Gobelin était juste derrière.
    Monsieur Je-Suis-Pas-Fatigué et madame Je-Veux-Pas-Dormir avaient décidé de rester en bas finalement.

    - Stop! Annonça monsieur Baba. C'est l'interdicte. Je suis le premier alors je suis le roi.

    Seul monsieur Je-Suis-Pas-Fatigué ne fut pas confus.
}
\imgpage{./img/page-05.jpg}

\txtpage{
    Monsieur Je-Suis-Pas-Fatigué hocha de la tête.

    - L'interdicte c'est parce que c'est l'édit d'arrêt proclamé par le roi de la tour, expliqua-t-il.

    Madame Gobelin n'était pas moins confuse.

    - Et c'est quoi un édit, demanda madame Gobelin?

    - Monsieur Je-Suis-Pas-Fatigué clarifia: Oh tu sais, c'est comme un dire mais on aspire fort
    avant, parce qu'il faut qu'il soit entendu de loin.
}
\imgpage{./img/page-06.jpg}

\txtpage{
    À l'heure du dîner, monsieur Je-Suis-Pas-Fatigué voulait qu'ils jouent de la musique,
    mais ils n'avaient pas encore réussi à s'entendre.
    Alors madame Gobelin et madame Je-Veux-Pas-Dormir étaient parties courir en attendant.
}
\imgpage{./img/page-07.jpg}

\txtpage{
    - Monsieur Baba proposa: Faisons un album. un album avant spectacle. Le nom? Guiguigoulax.

    - Monsieur Je-Suis-Pas-Fatigué réfléchit. C'est quoi Guiguigoulax?

    - Monsieur Baba, très sérieux, répondit: C'est un très vieux et mort.

    - Comme un hommage, demanda monsieur Je-Suis-Pas-Fatigué?

    - Monsieur Baba n'était pas du tout d'accord. C'est pas un hommage, C'EST UN ALBUM AVANT SPECTACLE.
}
\imgpage{./img/page-08.jpg}

\txtpage{
    Madame Gobelin avait entendu l'échange, et était revenue.

    - Elle demanda: Quel genre d'album?

    Monsieur Baba la regarda sérieusement

    - Une chanson c'est un haut parleur. C'est une lumière. C'est du son. C'est des batteries.

    Ça n'avait pas vraiment répondu à la question de madame Gobelin. En fait, elle soupçonnait quelque
    chose.

    - Dis, un album c'est quoi?
}
\imgpage{./img/page-09.jpg}

\txtpage{

    Monsieur Baba ouvrit la bouche pour répondre, mais monsieur Je-Suis-Pas-Fatigué était plus rapide.

    - Oui.

    Et ils jouèrent toutes sortes de musiques. Les passants étaient stupéfaits. Ou perplexes.
}
\imgpage{./img/page-10.jpg}

\txtpage{
    À présent, monsieur Baba voulu faire de la magie. Il sorti de sa poche un paquet de cartes,
    et se tourna vers monsieur Je-Suis-Pas-Fatigué.

    - Prends une carte.

    Mais monsieur Je-Suis-Pas-Fatigué était endormi!

    Monsieur Baba se tourna donc vers madame Je-Veux-Pas-Dormir, qui commençait à cogner des clous.

    - Donc je prends une carte, demanda-t-elle?

    Puis elle tira une carte.
}
\imgpage{./img/page-11.jpg}

\txtpage{
    Madame Je-Veux-Pas-Dormir révéla timidement sa carte.

    C'était le trois de carreaux!

    Madame Gobelin, qui était prête, se mit à entonner la chanson de la victoire.

    La chanson soudaine surpris madame Je-Veux-Pas-Dormir. Elle ne ressentait plus de fatigue.
    On pu aussi entendre monsieur Je-Suis-Pas-Fatigué, par terre, chantonner.

    Vraiment, le trois de carreaux est magique.
}
\imgpage{./img/page-12.jpg}

\txtpage{
    Pour s'étirer et se détendre, monsieur Baba et ses amis s'installèrent ensuite pour faire du
    yoga. Enfin, il fallut d'abord expliquer à monsieur Je-Suis-Pas-Fatigué ce que c'est car il
    jurait n'être pas fatigué.

    Monsieur Baba ferma ses yeux pour se concentrer.

    Avant qu'il puisse vraiment commencer, un drôle de bruit vint le brusquer.
}
\imgpage{./img/page-13.jpg}

\txtpage{
    - Beeeeeeeee, fit la drôle de créature.

    Monsieur Baba plissa des yeux pour essayer de comprendre.

    Il y avait... peut-être un petit mouton? C'était difficile à dire, il y avait tellement de
    pattes.
}
\imgpage{./img/page-14.jpg}

\txtpage{
    Monsieur Baba comptait les pattes: Un, deux, trois, ... mais il n'en finissait plus!

    Décidément, il devait bien y avoir... neuf milles pattes. Gros minimum.

    Il croyait au début que ce mouton était vraiment très touffu, mais en fait, il s'agissait même
    de pattes sur son dos! Ou... juste des bâtons.
}
\imgpage{./img/page-15.jpg}

\txtpage{
    - Monsieur Baba pris la parole: C'est quoi ton nom?

    - Beeeeeeee, répondit Matonrière.

    - Aaaah! Matonrière, répéta monsieur Baba.

    - Beeeeeee, fit encore Matonrière.

    - Tu veux que je monte sur ton dos? J'arrive.
}
\imgpage{./img/page-16.jpg}

\txtpage{
    Et il monta sur le dos de Matonrière. Elle se mit à courir vite, très vite. Ça court vite,
    neuf mille pattes. Monsieur Baba dû se cramponner aux bâtons pour ne pas tomber.

    Ils arrivèrent chez Matonrière, et elle fit à manger pour monsieur Baba. Elle ouvrit une canne
    de haricots, ils cassèrent la croûte, puis elle rangea les plats.

    - Beeeeeee, annonça Matonrière.

    Elle s'accroupit et...
}
\imgpage{./img/page-17.jpg}

\txtpage{
    Fit du yoga!
}

\legalpage{16 octobre 2026}
\backpage

\end{document}
